% Reconstructed text from V5 PDF page 30; non-native source
\par\smallskip\noindent\rule{\textwidth}{0.4pt}\par{\sffamily\footnotesize Original V5 page 30 -- recovered text}\par\begin{flushleft}\footnotesize\ttfamily
\par\smallskip
\par\smallskip
\noindent Diagnostics and Serving Infrastructure for MAS 30\par
\par\smallskip
\par\smallskip
\par\smallskip
\noindent 5.3 Cross-Model Cache Transfer\par
\noindent Cross-model cache transfer attempts to reuse KV-cache tensors computed by one model vari-\par
\noindent ant for inference on another, avoiding full recomputation when agents use different models or\par
\noindent adapter configurations.\par
\noindent  DroidSpeak [20] enables KV cache sharing between different fine-tuned variants of the same\par
\noindent base architecture. When transferring context from one LLM to another, DroidSpeak identifies\par
\noindent which transformer layers require fresh computation (layers where the KV representations di-\par
\noindent verge most between the two models) and which can reuse the donor model’s KV cache. This\par
\noindent selective layer-wise recomputation achieves up to 4× throughput improvement and 3.1× faster\par
\noindent prefill (TTFT reduction) with negligible quality loss on F1, Rouge-L, and code similarity met-\par
\noindent rics. The key requirement is architectural compatibility: the donor and recipient models must\par
\noindent share the same transformer architecture (different fine-tuned variants are acceptable, but cross-\par
\noindent architecture sharing is not supported). The system is implemented in vLLM (UChicago + Mi-\par
\noindent crosoft). For multi-agent scenarios, DroidSpeak enables agents using different LoRA-adapted\par
\noindent models to share cached context without full recomputation — a capability that could signifi-\par
\noindent cantly reduce serving costs in heterogeneous agent deployments.\par
\noindent  LRAgent (ICML 2026) manages KV cache for multi-adapter scenarios by separating a\par
\noindent shared base cache from per-adapter low-rank residuals, enabling cache reuse across requests\par
\noindent that share the same base model but differ in their LoRA adapters.\par
\par\smallskip
\noindent 5.4 Agent-Specific Extensions\par
\noindent Several systems extend general-purpose KV-cache management with agent-aware policies:\par
\par\smallskip
\noindent Agent-identity-aware scheduling. ScaleSim [28] manages heterogeneous per-agent state —\par
\noindent KV caches, LoRA adapters, auxiliary draft models for speculative decoding, and action his-\par
\noindent tories — using invocation-distance-based memory management. It is the only surveyed sys-\par
\noindent tem that manages non-KV agent state (adapters and auxiliary models) alongside KV cache\par
\noindent within a unified eviction framework. The core abstraction is invocation distance (§3.2): “a\par
\noindent unified abstraction that estimates the relative order in which agents will issue future LLM re-\par
\noindent quests.” Three categories of distance estimation are defined: independent simulation (inferred\par
\noindent from action duration), interaction-involved (D = min(Daction , Dinteraction ), where Dinteraction =\par
\noindent Physical Distance/Velocity), and predefined activation paths (hop count from information source\par
\noindent via BFS).\par
\noindent  The eviction policy is future reuse-aware: “instead of relying on recency-based policies\par
\noindent such as LRU, ScaleSim prioritizes eviction based on invocation distance” (§3.3). Agents with\par
\noindent larger distances are evicted first. A proactive prefetch mechanism restores offloaded agents be-\par
\noindent fore their activation, running in parallel with inference since LLM inference does not involve\par
\noindent CPU–GPU PCIe transfer. For shared memory objects (overlapping KV prefixes), the eviction\par
\noindent distance is set to the minimum across all referencing agents (Appendix A). A preemption fall-\par
\noindent back handles unexpected early activations by interrupting ongoing tasks.\par
\noindent  The system achieves up to 1.74× speedup over SGLang on information-diffusion simula-\par
\noindent tions, 1.73× on AgentSociety (independent agents), and 1.31× on Generative Agents (interaction-\par
\noindent involved), tested with Qwen2.5-7B on a single NVIDIA H100 at 20\% active agent rate (§4.2,\par
\noindent Figure 8). TTFT is consistently reduced by 48–68\% compared to hierarchical caching base-\par
\noindent lines. The speedup is bounded by the decode phase proportion of total time (§4.3). A key\par
\noindent limitation is that the application must provide invocation distance estimates through a unified\par
\end{flushleft}
